% !TeX root = ../main.tex

\sysusetup{
  keywords  = {视觉语言模型，多模态问答，智能助手，Qwen-VL-Plus，Gradio},
  keywords* = {Vision-Language Model, Multimodal QA, Intelligent Assistant, Qwen-VL-Plus, Gradio},
}

\begin{abstract}
  随着人工智能技术的快速发展，视觉语言模型（Vision-Language Model, VLM）在多模态理解领域取得了显著进展，为"看图问答"场景提供了可行的技术方案。本项目设计并实现了一个基于VLM的智能图文问答助手，支持用户上传图片（包括自然场景图像和文档截图）并进行中文多轮问答交互。系统采用阿里云DashScope API提供的Qwen-VL-Plus模型作为核心推理引擎，结合Gradio 6.x框架构建了ChatGPT风格的暗色主题Web交互界面，支持左侧会话列表管理与右侧聊天区域的双栏布局。项目实现了多轮对话、上下文记忆、图片识别与OCR、基础逻辑推理等核心功能，并通过构建包含真实自然场景和文档图像的80条评测数据集对系统进行了验证，整体准确率达到81.2\%。研究成果可应用于商品识别、讲义理解等实际场景，具有良好的应用前景。

\end{abstract}

\begin{abstract*}
  With the rapid development of artificial intelligence technology, Vision-Language Models (VLM) have made significant progress in multimodal understanding, providing a feasible technical solution for "image-based question answering" scenarios. This project designs and implements an intelligent image-text QA assistant based on VLM, supporting users to upload images (including natural scene images and document screenshots) and conduct Chinese multi-turn QA interactions. The system uses the Qwen-VL-Plus model provided by Alibaba Cloud DashScope API as the core inference engine, combined with the Gradio 6.x framework to build a ChatGPT-style dark-themed Web interface with a dual-panel layout featuring a session sidebar and a chat area. The project implements core functions such as multi-turn dialogue, context memory, image recognition and OCR, and basic logical reasoning, and verifies the system through an 80-item evaluation dataset containing real natural scene and document images, achieving an overall accuracy of 81.2\%. The research results can be applied to practical scenarios such as product recognition and lecture note understanding, with good application prospects.
\end{abstract*}